\section{Operational framework, vacuum waves, and null cuts}
\label{sec:foundations}

We use the metric $\eta=\operatorname{diag}(1,-1,-1,-1)$ and
$U(x)=e^{iP\cdot x}$. Scalar products are antilinear in their first
argument. Our dilation convention is
\[
 U_D(s)P_\mu U_D(-s)=e^sP_\mu,\qquad
 U_D(s)O(x)U_D(-s)=e^{sd}O(e^sx)
\]
for a field of dimension $d$. The distinction between a field, its
vacuum wave, and a bounded local operation will be maintained throughout.

\subsection{Original physical data and the vacuum-generated domain}

The hypotheses below specify the original fields, their bounded local
net, and their Ward currents. We also fix the domain required when a
new field is admitted to that theory. Proposition~\ref{F:old-domain}
constructs this domain from the original vacuum distributions and
proves its compatibility with the given affiliated realizations.

\Needspace{6\baselineskip}
\begin{assumption}[Operational physical theory]\label{ass:framework}
The physical Hilbert space $\cH$ is generated from a vacuum $\Omega$ by
the original physical fields and their adjoints. The fields, in finite-component
Lorentz multiplets, and their ordered products are operator-valued
Wightman distributions on a common invariant polynomial domain, which we
call the initial realization domain $\mathcal E$. They have positive translation spectrum,
Poincar\'e covariance, and locality,
with the usual grading when fermions occur. The vacuum is invariant under
Poincar\'e transformations and continuous dilations; uniqueness of the
invariant vector is not assumed.

Let $\mathcal D_{\mathrm{fact}}\subset\mathcal E$ be the algebraic span of
$\Omega$ and all finite original vacuum words with compactly supported
factor tests. To specify the vacuum-generated domain, for every finite ordered list $\mathbf A=(A_1,\ldots,A_n)$ of original fields and adjoints,
write $\Psi^{\mathrm{old}}_{\mathbf A}(h)$ for the continuous extension of
its original factor-test vacuum map to joint kernels
$h\in\mathcal S((\R^4)^n)$. These kernels
include $C_c^\infty((\R^4)^n)$ and finite sums of factorized tests. Set
\[
 \mathcal D_{\mathrm{old}}
 =\Span\{\Omega,\Psi^{\mathrm{old}}_{\mathbf A}(h):
          n=1,2,\ldots,\ h\in\mathcal S((\R^4)^n)\}.
\]
The existence of these extensions and their insertion actions is proved
in Proposition~\ref{F:old-domain}; it is not an additional assumption.
This is an algebraic span: it has no
Hilbert or graph completion and no bound on word length or tensor rank.
The domain whose containment is required below is precisely
$\mathcal D_{\mathrm{old}}$, not an independently larger choice of
$\mathcal E$.

Let $\mathcal F$ be the complex vector space of original point-field
labels, modulo equality of their operator-valued distributions on
$\mathcal D_{\mathrm{fact}}$; Proposition~\ref{F:old-domain} makes this
equivalently equality on $\mathcal D_{\mathrm{old}}$. It is the algebraic direct sum
\[
 \mathcal F=\bigoplus_{\Delta\in\Sigma}\mathcal F_\Delta,
 \qquad
 A=\sum_{\Delta\in S_A}A_\Delta,
 \quad S_A\subset\Sigma\ \text{finite},
\]
where $\Sigma\subset\R$ is a discrete set of real dimensions, bounded
below and locally finite: $\Sigma\cap[a,b]$ is finite for every finite
interval $[a,b]$. Each $\mathcal F_\Delta$ is a finite-dimensional
Lorentz-stable space of finite-component fields satisfying
\[
 U_D(s)A_\Delta(x)U_D(-s)
       =e^{s\Delta}A_\Delta(e^s x)
\]
as operator-valued distributions on $\mathcal D_{\mathrm{fact}}$, and hence
on $\mathcal D_{\mathrm{old}}$ by that proposition.
Multiplicity counts independent component fields after including all
Lorentz component indices and imposing linear field identities;
equivalently, each grade contains finitely many finite-component Lorentz
multiplets. Homogeneous components are themselves original physical
fields. Adjunction preserves each real grade, with conjugate Lorentz
components, and covariant distributional differentiation sends
$\mathcal F_\Delta$ into $\mathcal F_{\Delta+1}$, with a missing grade
interpreted as zero. Lorentz transformations preserve the grades and
statistics; translations act on spacetime arguments. No completion by
infinitely many scaling components is taken.

There is a fixed isotone net $O\mapsto\cM(O)$ of weakly closed
algebras of physical bounded operations, covariant under the same
physical Poincar\'e and dilation representations. Algebras at spacelike separation
commute, with the fixed statistics twist in the graded case. Each smeared
original physical field has a closed affiliated realization in its localization
algebra, extending its realization on $\mathcal E$, whose adjoint also
extends the field adjoint there. The
corresponding extensions on $\mathcal D_{\mathrm{old}}$ are derived in
Proposition~\ref{F:old-domain}.
For any other open region $R$,
\[
 \cM(R)=\bigvee_{\overline O\subset R}\cM(O),
\]
where $O$ ranges over double cones. Physical point-field membership in
$\Aphys$ is intrinsic to this fixed net, not to an enumeration of symbols.
An operationally reconstructed finite family of homogeneous
finite-component point fields of real scaling dimensions is admitted,
together with its finite linear span, once its localization supports can
shrink inside arbitrarily small double cones; all joint field distributions
of the family with itself and with every fixed finite list of original
fields and adjoints have been established at every finite length on a
compatible common invariant domain containing $\mathcal D_{\mathrm{old}}$,
with the required adjoints and the preceding covariance, spectrum, and
locality/statistics properties; and compatible smeared realizations
affiliated with the original localization algebras have been established,
with their adjoints extending the field adjoints on that domain. Separately
constructed fields are not combined into a common polynomial field algebra
without their mutual joint distributions, common invariant domain with the
required adjoints, and original-net affiliations having been established.
The criterion asserts no candidate's existence or analytic properties and
admits no point field merely by convergence of an infinite nonhomogeneous
scaling sum.

The original physical fields include a symmetric conserved stress tensor
$T_{\mu\nu}$ with the translation and orbital Lorentz Ward identities,
and a local virial field $V_\mu$ such that
\[
 D_\mu=x^\nu T_{\mu\nu}-V_\mu,\qquad
 \partial^\mu D_\mu=0,\qquad
 T^\mu{}_{\mu}=\partial^\mu V_\mu.
\]
The Ward identities of $D$ generate the stated dilations. Charge Ward
identities are understood through compact Cauchy slabs surrounding an
insertion and extend to polynomial products by the derivation rule.
The identities displayed above are flat-space field identities; the
usual contact terms in time-ordered distributions are allowed.
\end{assumption}

Finite scaling support concerns field labels. It leaves all finite product
lengths and all joint Schwartz and compact smooth kernels available,
without a finite tensor-rank requirement. It restricts neither Hilbert
completions and generator spectra nor weak closures of bounded algebras.
The bounded-locality and original-net affiliation clauses are separate
hypotheses: neither weak commutation of field restrictions nor affiliation
in a larger relatively local net supplies them. No Haag duality, split
property, phase-space estimate, scattering theory, or conformal symmetry
is assumed.

\Needspace{6\baselineskip}
\begin{proposition}[Original joint-test domain]\label{F:old-domain}
For every fixed finite ordered list of original fields and adjoints,
the vacuum distribution has a unique tempered extension and its vacuum
map has a continuous Hilbert-valued extension to all joint Schwartz
kernels. The space $\mathcal D_{\mathrm{old}}$ is Hilbert-dense and
invariant under old field insertions and continuous joint polynomial
insertions, with their prescribed adjoints. Finite field-label
identities, covariance, locality, and the compact-slab Ward identities
pass to this domain with their stated support conditions.

For a compactly smeared original field $B(f)$, each retained closed
affiliated realization and its Hilbert adjoint agree on
$\mathcal D_{\mathrm{old}}$ with these insertion actions. This assertion
also holds for an originally supplied Schwartz realization. It does
not require $\mathcal D_{\mathrm{old}}$ to be a graph core for that
realization.
\end{proposition}

\begin{proof}
\emph{Tempered distributions and finite field-label expansions.}
For every fixed ordered list of homogeneous original fields
$A_i\in\mathcal F_{\Delta_i}$, its already given vacuum distribution is
tempered. To see this, remove the common translation coordinate. In the
remaining relative coordinates $\xi$, dilation invariance of the vacuum
gives
\[
 W(\lambda\xi)=\lambda^{-\sum_i\Delta_i}W(\xi),\qquad \lambda>0,
\]
in the distributional pullback sense. A compact neighborhood of the
origin is controlled by the local finite order of the given distribution.
On a fixed compact annulus it is bounded by finitely many test-function
seminorms. Dilating that annulus transfers the bound to every outer dyadic
annulus with a fixed power determined by the dimensions and the seminorm
order. Summation against a sufficiently high Schwartz seminorm proves
temperedness. The zero- and one-point cases are immediate.

Adjoints are homogeneous at the same real grade. For any fixed list of
original fields $A_i=\sum_{\Delta\in S_i}A_{i,\Delta}$, multilinearity on
$\mathcal D_{\mathrm{fact}}$ first gives the component expansion on compact
factor tests. Compact tensor-test density and continuity of the given joint
distributions then give, for the same compact joint kernel $\varphi$,
\begin{equation}\label{F:finite-label-extension}
 W(A_1,\ldots,A_n)[\varphi]
 =\sum_{(\Delta_1,\ldots,\Delta_n)\in S_1\times\cdots\times S_n}
 W(A_{1,\Delta_1},\ldots,A_{n,\Delta_n})[\varphi].
\end{equation}
The finite sum of homogeneous tempered extensions defines the unique
tempered extension of the left side, so the same identity holds for every
joint Schwartz kernel. It agrees with any originally specified Schwartz
distribution by its stated continuity and compact-test density. Its constants
and seminorm orders may depend on the fixed list; no single total scaling
dimension is assigned to a nonhomogeneous sum.

\emph{Hilbert-valued words and insertion actions.}
For a fixed old list $\mathbf A$, write $\Psi^0_{\mathbf A}(u)$ for the
original vacuum vector of a finite sum $u$ of compact factor tests. The
doubled tempered distribution gives a continuous Schwartz seminorm $p$
and a constant $C$ such that
\[
 \|\Psi^0_{\mathbf A}(u)-\Psi^0_{\mathbf A}(v)\|^2
 =W(\mathbf A^*,\mathbf A)[(u-v)^*\otimes(u-v)]
 \le C p(u-v)^2.
\]
Here the star reverses and adjoints the field list and conjugates and
reverses the kernel. Compact finite tensor tests are dense in the joint
Schwartz space. The estimate therefore constructs the unique continuous
Hilbert-valued extension $\Psi^{\mathrm{old}}_{\mathbf A}$ used in the
definition of $\mathcal D_{\mathrm{old}}$; polarization passes all its Gram
identities to the limit. Any originally specified Schwartz vacuum-word map
agrees with this extension by its stated continuity and the same density.
Consequently the original vacuum-generation assumption also makes
$\mathcal D_{\mathrm{fact}}$ Hilbert-dense.
This argument concerns only original products whose existence is assumed;
it supplies no product containing a newly constructed field.

For a fixed old word
$\mathbf A$ and a fixed finite old prefix $\mathbf B$, the same doubled
vacuum-distribution estimate applies to the longer word. Consequently, if
$h_j\to h$ in the joint Schwartz topology, then for a fixed prefix kernel
$g$,
\[
 \Psi^{\mathrm{old}}_{\mathbf A}(h_j)\longrightarrow
 \Psi^{\mathrm{old}}_{\mathbf A}(h),\qquad
 \Psi^{\mathrm{old}}_{\mathbf B,\mathbf A}(g\otimes h_j)
 \longrightarrow
 \Psi^{\mathrm{old}}_{\mathbf B,\mathbf A}(g\otimes h)
\]
in Hilbert norm. The estimates are continuous jointly in $g,h$ on their
Schwartz spaces. Reversed adjoint kernels obey the same bounds. Compact
joint kernels can be approximated, after a finite partition into product
patches, by finite sums of factor tests with the prescribed local supports;
Schwartz kernels admit finite tensor approximations in the Schwartz
topology. Concatenation, reversal, and differentiation are continuous in
these test spaces.

Left concatenation of old test words is well-defined on their vacuum
vectors. Indeed, if a finite old test polynomial $P$ represents zero,
positivity and Cauchy--Schwarz applied to $P$ and $Q^*QP$ give
$W(P^*Q^*QP)=0$ for every finite old test polynomial $Q$. This is the
squared norm of the proposed vector for $QP$. Thus the null vectors form
a left ideal, including cancellations between different words, and the
same vacuum pairing places adjoint insertion inside the Hilbert adjoint.
Concatenation therefore makes $\mathcal D_{\mathrm{old}}$ invariant
under the old field and polynomial insertions. It is Hilbert-dense because
it contains the original factorized vacuum-polynomial vectors that generate
$\cH$.

\emph{Agreement with the retained realizations.}
Let $T_B(f)$
be a chosen closed affiliated realization of an original field with compact
test $f$, and choose finite sums of compact tensor tests $h_j\to h$ as above.
On the vectors in $\mathcal D_{\mathrm{fact}}\subset\mathcal E$ the original extension clause gives
\[
 T_B(f)\Psi^{\mathrm{old}}_{\mathbf A}(h_j)
   =\Psi^{\mathrm{old}}_{B,\mathbf A}(f\otimes h_j).
\]
Both sides converge in Hilbert norm by the two displayed continuity
statements. Closedness of $T_B(f)$ therefore proves the same equality for
$h$. Applying the argument to the closed Hilbert adjoint $T_B(f)^*$ gives
its equality with the insertion of $B^\dagger(\bar f)$ there. Iterating
this one-field statement proves the corresponding actions of every finite
factorized prefix and its reversed adjoint on
$\mathcal D_{\mathrm{old}}$. Joint polynomial insertions are the continuous
concatenations already defined; this argument does not assign an affiliated
closure to an arbitrary joint smear.

If an original Schwartz smearing $B(f)$ and its closed affiliated
realization were already supplied, the same conclusion holds for that
realization. To identify its action on a fixed vector of
$\mathcal D_{\mathrm{fact}}$, approximate $f$ by compact tests in Schwartz
topology. The longer doubled correlations give Hilbert convergence of the
inserted vectors, while the original Schwartz distributional continuity
gives the same limiting matrix elements against
$\mathcal D_{\mathrm{fact}}$. Density identifies the vectors, also for the
adjoint. The preceding $h_j$ closed-graph argument then applies unchanged.
This uses only an originally supplied Schwartz action and its stated
continuity; it asserts no pre-existing realization for a smearing that was
not originally supplied.

\emph{Identities and local Ward actions.}
The same paired vector and inserted-vector limits extend every finite
field-label identity from $\mathcal D_{\mathrm{fact}}$ to
$\mathcal D_{\mathrm{old}}$. The converse follows from
$\mathcal D_{\mathrm{fact}}\subset\mathcal D_{\mathrm{old}}$, so the two
equality conventions for $\mathcal F$ agree. The homogeneous dilation law
and the adjoint and derivative grade laws extend in the same way, because
the transformed, reversed, and differentiated kernels remain in the stated
test spaces. Thus none of these identities on joint vectors was used to
construct those vectors.

The original covariance and flat field identities pass from the factorized
vectors to joint vectors by the same longer-word continuity, with transformed
or differentiated kernels. Locality and compact-slab Ward identities pass
under their stated support conditions. For a general Schwartz spectator
vector, keep the localized slab and insertion tests fixed and approximate
only the spectator kernel in Schwartz topology. If a localized insertion
polynomial itself has a compact joint kernel, use the support-preserving
compact approximants.
Thus the operational identities used below hold on the dense invariant
vacuum-generated domain $\mathcal D_{\mathrm{old}}$, with the original
vacuum correlations, Hilbert space, and bounded net unchanged. The
factor approximations just used converge in the graph of each specified
$T_B(f)$ on these joint vectors; Hilbert density alone was not used for
that conclusion. No claim is made that $\mathcal D_{\mathrm{old}}$ is a
graph core for a chosen realization, that the closure of its field
restriction equals that realization, or that a newly constructed field acts on arbitrary additional
vectors of $\mathcal E$.
\end{proof}

We will use \eqref{F:finite-label-extension} as a finite-label permanence
principle for later distributions. Every term uses the same joint kernel;
continuity, common spectral cones, and linear distributional identities
pass through the finite sum. Positivity and descent through null
polynomial vectors require their full summed forms and will be checked
where they arise.

The dimension set is countable, and each grade has a finite component
basis. Countable dense test families for the resulting countable
collection of finite words, together with the Hilbert-valued continuity
just proved, show that $\cH$ is separable. This supplies the separability
used in the direct-integral descriptions below.

The spectrum condition now gives tube analyticity, uniqueness from an open
boundary patch, and Reeh--Schlieder for these original distributions. In
particular, a separated-point field identity valid on a mutually spacelike
open patch extends to ordered vacuum distributions by tube uniqueness.
This assertion does not assign differentiated Ward identities to a
particular choice of time ordering.

\subsection{Example: the full free scalar theory}
\label{F:free-scalar-example}

The massless real scalar in its four-dimensional vacuum Fock
representation gives a concrete example. Take the full scalar net
$\cM(O)=\{e^{i\overline{\phi(f)}}:
f\in C_c^\infty(O,\R)\}^{\prime\prime}$ and, as original fields, finite
linear combinations of Wick monomials
$:\!\prod_{j=1}^n\partial^{\alpha_j}\phi\!:$.
Their dimensions are the nonnegative integers
$n+\sum_j|\alpha_j|$, with finitely many component monomials at
each grade. The Weyl relations give bounded locality, and Wick's
construction supplies all joint products and their invariant domain
\cite[Theorem~5.7 and its following remark]{BrunettiFredenhagenKoehler1996}.
Affiliation has a concrete justification here. On the common domain
invariant under Wick fields and Weyl operators
\cite[Section~2.2, equation~(11)]{BrunettiFredenhagen2000},
a Wick smearing commutes with exterior Weyl operators. Its closed
graph retains this commutation and domain invariance. Double-cone
duality of the free massless scalar
\cite[Section~3, Corollary~2]{HislopLongo1982} identifies the generated
exterior algebra with $\cM(O)'$, proving affiliation. The formal
adjoint gives the required Hilbert-adjoint extension. This is a
property of the example, not a duality hypothesis of the theorem.

\subsection{Selecting the current weights}
\label{F:weights}

We may, and henceforth do, replace $T$ and $V$ by Hermitian homogeneous
fields of dimensions four and three. We justify both the selection and
preservation of the charges. First take the even components of the
spacetime currents. The prescribed unitary affine derivations preserve
parity, so the odd-current Ward contributions vanish by parity separation
on parity-homogeneous insertions; the flat field identities split by parity as
well. For an even current $J$ with a real slab test, polynomial adjunction
on $\mathcal D_{\mathrm{old}}$ gives
\[
 \bigl(i[J,A]\bigr)^\dagger=i[J^\dagger,A^\dagger].
\]
The prescribed unitary derivation preserves adjoints, and its Ward identity
also applies to $A^\dagger$ with the same support. Hence $J^\dagger$ and
$(J+J^\dagger)/2$ have the same Ward action as $J$. Applying this to the
translation and orbital currents of $T$ and to $D=xT-V$ shows that the
Hermitian parts preserve all required charges. The real metric, coordinate,
and derivative coefficients also preserve the flat field identities under
adjunction. Relabel these even Hermitian representatives as $T,V$.
The adjoint grade law preserves their decompositions. Write the finite sums
\[
 T=\sum_{d\in S_T}T^{(d)},\qquad
 V=\sum_{e\in S_V}V^{(e)},
\]
with every missing component interpreted as zero. Directness of the grading,
Lorentz covariance, and the derivative grade law split symmetry,
Hermiticity, conservation, and the trace identity. Thus, for every
$d\in S_T\cup\{e+1:e\in S_V\}$,
\[
 T^{(d)}_{\mu\nu}=T^{(d)}_{\nu\mu},\qquad
 \partial^\mu T^{(d)}_{\mu\nu}=0,\qquad
 (T^{(d)})^\mu{}_{\mu}=\partial^\mu V^{(d-1)}_\mu,
\]
and the components can be taken Hermitian. In particular each geometric
current
\[
 D^{(d)}_\mu=x^\nu T^{(d)}_{\mu\nu}-V^{(d-1)}_\mu
\]
is conserved.

Fix a homogeneous insertion $O$ of dimension $\alpha$ and homogeneous
spectators. A compact-slab translation commutator made from $T^{(d)}$ has
weight $\alpha+d-3$: the slab integral removes three powers of length.
An orbital Lorentz commutator from $xT^{(d)}$ has weight $\alpha+d-4$,
as does a dilation commutator from the pair $D^{(d)}$. Conservation of
these component currents permits the slabs to be rescaled while keeping
the insertion inside. After testing the homogeneous spectators, these are
identities between ordinary scalar distributions. The prescribed
translation, Lorentz, and dilation actions have weights $\alpha+1$,
$\alpha$, and $\alpha$, respectively. After testing the remaining variables,
the Ward identities are finite scalar sums of real exponentials in the
logarithmic scale. Their linear independence
selects the $d=4$ translation and orbital contributions and the $d=4$
dilation pair $xT^{(4)}-V^{(3)}$; the other corresponding Ward
contributions vanish. This is a statement about charges, not the vanishing
of the nonselected fields.

Finite multilinearity now extends these identities to every fixed list of
original fields and adjoints under the same test kernels, and the retained
derivation rule extends them to every finite polynomial word. The selected
components retain conservation, the trace identity, and all required Ward
actions. Replacing the original stress tensor by $T^{(4)}$ adds a local
symmetric conserved tensor with zero translation and Lorentz charges.
This is an admissible charge-preserving change of stress tensor; it does
not presume a scalar improvement. The field decompositions here are
identities on $\mathcal D_{\mathrm{old}}$; no equality between a chosen
closed realization of a sum and a sum of chosen closed realizations is
needed. Each selected component has its own affiliated realization by
Assumption~\ref{ass:framework}.

\subsection{Modular input and a spectral comparison lemma}
\label{F:modular}

For a standard real subspace $X\subset\cH$, write
\[
 \mathfrak s_X(x+iy)=x-iy\quad(x,y\in X),\qquad
 \mathfrak s_X=J_X\Delta_X^{1/2},\qquad
 S_X=-\frac{1}{2\pi}\log\Delta_X.
\]
Here standard means $\overline{X+iX}=\cH$ and $X\cap iX=\{0\}$.
The symplectic complement is
$X'=\{\xi:\operatorname{Im}\ip{\xi}{x}=0\text{ for }x\in X\}$.
For a standard algebra $(\cM,\Omega)$ its associated real space is
$X_\cM=\overline{\{A\Omega:A=A^*\in\cM\}}$.
Statistics twists in complementary spaces and algebras are understood
in all statements involving odd fields.

Every original polynomial word in this subsection has the finite component
expansion described above, with its joint test kernel unchanged. Thus the
preceding temperedness and finite-component covariance apply to every word
used in the polynomial argument.

The Wightman Bisognano--Wichmann theorem identifies the wedge polynomial
Tomita graph with $J_W e^{-\pi B}$, where $B$ generates the wedge boost;
hence $\Delta_W^{it}=e^{-2\pi itB}$ and the opposite wedge has the
complementary real space \cite{BW1976}. The published theorem assumes a
unique translation-invariant vacuum. We explain its boundary calculation
and core argument in the version needed here, followed by the passage
to bounded algebras. First, uniqueness of the vacuum is unnecessary for
this graph identification. The reflection itself must be constructed before the
core argument. Extend the covariance of the vacuum distributions along
the connected complex Lorentz group. On a real Jost patch, locality
reverses the order of the fields with the prescribed statistics signs;
analytic uniqueness extends the resulting distributional equality.
Combine it with the Hermitian-conjugation identity for vacuum
distributions. If $F$ and $G$ are smeared polynomial words and
$\xi_F=F\Omega$, it says
\[
 \ip{\xi_F^{\natural}}{\xi_G^{\natural}}
       =\ip{\xi_G}{\xi_F}.
\]
Here $\xi_F^{\natural}$ is the reflected-adjoint word: conjugate its
test functions, reflect their arguments by
$j_W(u,v,y)=(-u,-v,y)$, and replace the components by their adjoints
with the finite-component matrices of the imaginary wedge boost at
rapidity $i\pi$ and the fixed statistics twist. Apply this operation
to each factor in the \emph{original word order}; the order reversal
above occurs in the vacuum distributions used to compute scalar
products. This is precisely the finite-component reflection boundary
calculation in the polynomial BW proof. It uses covariance, the spectrum
condition, locality, and Hermitian conjugation, with no cluster limit
or condition on the dimension of the invariant subspace.

The displayed identity shows that the conjugate-linear word map
preserves norms and descends through all zero-norm polynomial relations.
Reflected test functions again range over the full test space, so its
range is dense. It therefore extends to an antiunitary $J_W$;
reflection twice, with the normal-statistics phases, gives $J_W^2=1$.
Now use polynomial vacuum vectors whose first point and successive gaps
lie in an open wedge-ordered patch. Imaginary boost
through rapidity $\pi$ takes their ordered distributions through the
positive-spectrum tubes, and locality at the spacelike boundary reverses
the fields. Tube uniqueness makes their complex span Hilbert-dense.
Adjoin the adjoint vectors and all real boost translates. The resulting
span is boost-invariant and its adjunction agrees with
$J_W e^{-\pi B}$. To check that it is a graph core, let $\mathcal G$ be
its closed graph and smooth by $e^{-\epsilon B^2}$. This smoothing
preserves $\mathcal G$, and both
$e^{-\epsilon B^2}$ and $e^{-\pi B}e^{-\epsilon B^2}$ are bounded.
Hilbert approximation by the indicated span consequently approximates
every smoothed vector in graph norm. Finally
$e^{-\epsilon B^2}\xi\to\xi$ in the graph norm of $e^{-\pi B}$ for
$\xi$ in its domain. This proves graph totality without a clustering
argument.

Second, the polynomial and bounded-algebra graphs coincide under
Assumption~\ref{ass:framework}. Start with finite sums of factorized
smeared words. If $A$ is a closed affiliated realization of such a
polynomial with $\Omega\in\Dom A\cap\Dom A^*$, its bounded polar
cutoffs $A_m=A\mathbf1_{[0,m]}(|A|)$ satisfy
\[
 A_m\Omega\longrightarrow A\Omega,\qquad
 A_m^*\Omega\longrightarrow A^*\Omega.
\]
For the second convergence use
$A_m^*\Omega=\mathbf1_{[0,m]}(|A|)A^*\Omega$.
Use natural-domain products and sums of the chosen closed affiliated
single-field realizations. Their domains contain $\mathcal D_{\mathrm{old}}$,
and their adjoints contain the corresponding adjoint polynomials there,
so they are closable. Their natural domains and graphs are invariant
under commutant unitaries; the closures remain affiliated. These are
the extensions to which the same polar-cutoff argument applies.
For an admitted compact joint kernel, first use a finite partition into
product patches inside the localization region, then approximate each part
by finite sums of tensor tests in its smooth test-function topology. The
reversed adjoint kernels of these approximants converge to the reversed
adjoint kernel in the same topology. The old Hilbert-valued
joint-distribution bounds proved above pass both the vacuum and
adjoint-vacuum vectors to the limit. Closedness of the bounded-algebra
Tomita graph therefore gives the same graph inclusion for the joint word.
This vector-level limit does not assert a closed affiliated realization for
an arbitrary joint smear.
Let $X_{\rm pol}(R)$ be the closed real span of vacuum vectors of finite
original polynomial expressions that are Hermitian on
$\mathcal D_{\mathrm{old}}$ and compactly localized in $R$, with
admitted compact joint kernels included by the graph limits just proved.
This is an auxiliary vacuum-word graph space; it neither restricts the
original test-function spaces nor redefines $\mathcal D_{\mathrm{old}}$.
Consequently,
\begin{equation}\label{F:BW-sandwich}
 X_{\rm pol}(W)\subset X_{\cM(W)}
 \subset X_{\cM(W')}^{\prime}
 \subset X_{\rm pol}(W')^{\prime}=X_{\rm pol}(W).
\end{equation}
This proves equality and the bounded wedge modular formula. The
inclusion in the middle is bounded locality, not a claim about commuting
unbounded restrictions. To pass from graph equality to algebra wedge
duality, use the modular-invariant-subalgebra theorem: a subalgebra
invariant under the modular group of a faithful normal state admits a
state-preserving conditional expectation \cite{Takesaki1972}; see also
\cite[Theorems~3.5 and~5.2]{AccardiCecchini1982}. If the subalgebra is
$\mathcal N$ and $E$ is this expectation, its GNS implementation is the
vector identity
\[
 E(x)\Omega=P_{\overline{\mathcal N\Omega}}x\Omega.
\]
Apply it to the local inclusion of the
opposite, twist-corrected wedge algebra in $\cM(W)'$. The larger
algebra's modular group is the geometric opposite boost and preserves
the smaller algebra. The latter has cyclic vacuum, so the implementing
projection is the identity. The expectation is therefore the identity
by separation, proving equality of these algebras. Double-cone duality
is not used.

We use the following exact forms of the modular translation theorems
\cite{Borchers1992,ArakiZsido2005}; for the positive-translation
statement we use the explicit formulation and proof in
\cite[Theorem~II.5.2]{Borchers2000}. If $e^{icb}$ fixes $\Omega$, $b\ge0$,
and $\Ad(e^{icb})\cM\subset\cM$ for $c\ge0$, where $\Omega$ is cyclic
and separating for $\cM$, then
\begin{equation}\label{F:Borchers}
 e^{isS_\cM}b e^{-isS_\cM}=e^s b,\qquad J_\cM bJ_\cM=b.
\end{equation}
If $\mathcal N\subset\cM$ have the same cyclic separating vector and
$\Ad(e^{isS_\cM})\mathcal N\subset\mathcal N$ for $s\ge0$, the
half-sided modular inclusion theorem gives
\begin{equation}\label{F:HSMI}
 \begin{split}
 H&=\overline{S_\cM-S_\mathcal N}\ge0,\qquad H\Omega=0,\qquad
 e^{isS_\cM}He^{-isS_\cM}=e^sH,\\
 J_\cM HJ_\cM&=H,\qquad
 \mathcal N=\Ad(e^{iH})\cM,\qquad
 e^{isS_\mathcal N}=e^{i(1-e^s)H}e^{isS_\cM}.
 \end{split}
\end{equation}
The closure in the first line is the self-adjoint closure on the common
logarithm domain; subsequent sums of affine-group generators will have
this meaning. The applicable inclusion result is
\cite[Theorem~2.1]{ArakiZsido2005}, specialized to a vector state. It
requires no cyclicity for $\mathcal N'\cap\cM$. Since $J_\cM$ is
antiunitary, its unitary-group formula is
$J_\cM e^{itH}J_\cM=e^{-itH}$, consistently with the generator formula
above. These results also hold for standard real subspaces; one may
apply the algebra result to their CCR Weyl algebras and restrict the
vacuum modular objects to the one-particle space.

For later form comparisons, $X\subset X_1$ implies
$\mathfrak s_X\subset\mathfrak s_{X_1}$. Thus the closed form of
$\Delta_X$ restricts that of $\Delta_{X_1}$ and
$\Delta_X\ge\Delta_{X_1}$ in form order. The resolvent integral for the
operator-monotone logarithm gives $S_X\le S_{X_1}$ on common operator
domains. This statement by itself is weaker than the next lemma. Koot's
spectral characterization of inclusions of standard pairs uses the
modular-strip and half-plane mechanism
\cite[Theorem~3.1, Lemma~3.4 and Theorem~3.5]{Koot2025}.
We prove the contraction version between two possibly different
standard spaces that is needed for the embedded comparison channels.

\Needspace{6\baselineskip}
\begin{lemma}[Strip comparison]\label{F:strip-comparison}
Let $X_i$ be standard real spaces with normalized modular generators
$B_i$, and let $K_i\ge0$ be self-adjoint with
$e^{itB_i}K_i e^{-itB_i}=e^tK_i$. Suppose $j:\cH_1\to\cH_2$ is a
contraction intertwining the boost groups and
\[
 j e^{iK_1}X_1\subset e^{iK_2}X_2.
\]
Then, for every $R\ge0$,
\begin{equation}\label{F:strip-bands}
 \mathbf1_{(R,\infty)}(K_2)j\mathbf1_{[0,R]}(K_1)=0.
\end{equation}
In particular, for $s>0$ and $\xi\in\Dom K_1^s$,
$j\xi\in\Dom K_2^s$ and
$\norm{K_2^s j\xi}\le\norm{K_1^s\xi}$.
\end{lemma}

\begin{proof}
We first prove that, for every $r>0$,
\begin{equation}\label{F:strip-exponential}
 e^{rK_2}j e^{-rK_1}\ \text{is everywhere defined and contractive}.
\end{equation}
Its exponential growth on a high target band will then exclude that band
from the image of a low source band.

Set $A=e^{-iK_2}je^{iK_1}$. It maps $X_1$ into $X_2$, so
$\mathfrak s_2A\xi=A\mathfrak s_1\xi$ on
$\Dom\mathfrak s_1=\Dom\Delta_1^{1/2}$. Equivalently,
\[
 \Delta_2^{1/2}A\xi
 =J_2AJ_1\Delta_1^{1/2}\xi
 \qquad(\xi\in\Dom\Delta_1^{1/2}).
\]
Thus $\Delta_2^{1/2}A\Delta_1^{-1/2}$, initially on
$\Dom\Delta_1^{-1/2}$, has the contractive bounded extension
$J_2AJ_1$. The operators
\[
 F(z)=\Delta_2^{iz}A\Delta_1^{-iz},\qquad
 -\tfrac12<\operatorname{Im}z<0,
\]
have contractive boundary values: the real boundary multiplies $A$
on either side by a unitary, and the other boundary is
$\Delta_2^{it}J_2AJ_1\Delta_1^{-it}$. The analytic-extension theorem
\cite[Theorem~2.3]{ArakiZsido2005}, followed by three lines, gives the
contractive analytic continuation to the strip. Equivalently, sandwich
between bounded spectral intervals of the modular logarithms, apply
three lines to their scalar matrix coefficients, and remove the
intervals. The displayed Tomita relation supplies the full source domain
needed at the second boundary.
Write $\widetilde F(\zeta)=F(z)$ for $\zeta=e^{-2\pi z}$, which
maps the strip bijectively onto the upper half-plane. Boost covariance
and the intertwining property of $j$ identify its positive real boundary
value as $e^{-i\zeta K_2}j e^{i\zeta K_1}$. To justify the interior
expression despite its potentially unbounded first factor, take
$\xi$ and $\eta$ in bounded spectral bands of $K_1$ and $K_2$,
respectively. The scalar function
\[
 \ip{e^{i\overline\zeta K_2}\eta}{j e^{i\zeta K_1}\xi}
\]
is entire and agrees with
$\ip{\eta}{\widetilde F(\zeta)\xi}$ on a positive real interval.
Boundary uniqueness, or reflection of their zero difference across
that interval, identifies them throughout the upper half-plane. At
$\zeta=ir$, $r>0$, contractivity therefore gives
\[
 \left|\ip{e^{rK_2}\eta}{j e^{-rK_1}\xi}\right|
 \le\norm{\eta}\,\norm{\xi}.
\]
Bounded spectral vectors are a core for $e^{rK_2}$. The adjoint-domain
criterion implies $j e^{-rK_1}\xi\in\Dom e^{rK_2}$ and identifies its
image with $\widetilde F(ir)\xi$. Finally approximate any source vector
by its $K_1$ spectral truncations. Boundedness of $e^{-rK_1}$ and
$\widetilde F(ir)$, followed by closedness of $e^{rK_2}$, proves that
$e^{rK_2}je^{-rK_1}$ is everywhere defined and contractive.
If $\xi=\mathbf1_{[0,R]}(K_1)\xi$, apply this extension to
$e^{rK_1}\xi$ to obtain
\[
 \norm{e^{rK_2}j\xi}\le e^{rR}\norm{\xi}.
\]
Projection onto $K_2\ge R+\epsilon$ and $r\to\infty$ shows that this
projection annihilates $j\xi$. Taking $\epsilon\downarrow0$ proves
\eqref{F:strip-bands}. The resulting tail inequality is
\[
 \norm{\mathbf1_{(a,\infty)}(K_2)j\xi}
 \le\norm{\mathbf1_{(a,\infty)}(K_1)\xi}.
\]
Integrate its square against $2s a^{2s-1}\dd a$. Truncation followed by
monotone convergence proves both the domain assertion and the norm
bound.
\end{proof}

\subsection{Timelike vacuum channels}
\label{F:waves}

Fix future null vectors $n,\bar n$ with $n\cdot\bar n=1$ and orthonormal
transverse axes $e_i$, $i=1,2$, with $e_i\cdot e_j=-\delta_{ij}$. Write
\begin{equation}\label{F:null-coordinates}
 \begin{gathered}
 x=vn+u\bar n+y,\qquad p=an+b\bar n+k,\qquad
 p\cdot x=au+bv-\boldsymbol k\cdot\boldsymbol y,\\
 2ab=M^2+|\boldsymbol k|^2.
 \end{gathered}
\end{equation}
The variables $a,b,\boldsymbol k$ also denote their commuting momentum
operators. Coordinate subscripts on covariant tensors mean contraction
with the coordinate basis vectors. The boost $B$ acts on spacetime by
$v\partial_v-u\partial_u$, and null rotations $G_i$ by
$u\partial_i+y_i\partial_v$.

Disintegrate the positive two-point measure of a dimension-$d$ multiplet
on the open timelike cone. Lorentz transformations and dilations act
transitively there, so relative to $\dd^4p$ its vacuum amplitude at rest
has norm proportional to $M^{d-2}$. The proportionality is a positive
finite-dimensional rotation-invariant matrix; transport to other
momenta uses the finite Lorentz matrices and unitary fiber transport.
There is no independent positive singular measure on an interior mass
shell, since dilation moves every positive mass shell and fixes the
homogeneous measure class on their union.

A possible contribution on the nonzero light cone is normalized instead
by $\dd^4p\,\theta(p^0)\delta(p^2)$. Its amplitude on a fixed null ray
would grow as energy to the power $d-1$. A boost along that ray has only
the boost weights of the finite Lorentz multiplet; positivity therefore
excludes the contribution if $d-1$ exceeds their maximum. The maxima
are one for a vector and two for a symmetric rank-two tensor, whereas
$d-1$ is respectively two and three for $V$ and $T$. A zero-momentum
atom for a positive-dimensional field is also impossible: its norm is
invariant under unitary dilation but would be multiplied by $e^{sd}$.
Thus both $V\Omega$ and $T\Omega$ have purely timelike spectrum.

Let $\Theta=T^\mu{}_{\mu}$. At this stage define only vacuum waves,
\begin{equation}\label{F:wave-decomposition}
 l=-\Theta/M^2,\qquad
 J_\mu=V_\mu-ip_\mu l,\qquad
 t_{\mu\nu}=T_{\mu\nu}
       -\tfrac13(-\eta_{\mu\nu}M^2+p_\mu p_\nu)l.
\end{equation}
The local Ward identities imply
$p^\mu J_\mu=0$, $p^\mu t_{\mu\nu}=0$, and
$\eta^{\mu\nu}t_{\mu\nu}=0$. At rest their rotation types are respectively
spin zero, one, and two, so they are orthogonal. Their rest norms are
constant, proportional to $M$, and proportional to $M^2$. In particular
$l$ is a well-defined dimension-two wave; division by $M^2$ has introduced
no infrared singularity in its two-point measure.

These channels have the two-point data of a scalar primary of dimension
two, a conserved vector primary of dimension three, and a conserved
symmetric traceless primary of dimension four. This follows directly
from the uniqueness of a positive rotation-invariant metric on each of
the three rest representations and homogeneous Lorentz transport.
Equivalently their position-space kernels are the positive-energy
boundary values proportional to
\[
 \frac1{(x^2)^2},\qquad
 \frac{I_{\mu\nu}(x)}{(x^2)^3},\qquad
 \frac{I_{\mu\rho}I_{\nu\sigma}+I_{\mu\sigma}I_{\nu\rho}
                 -\eta_{\mu\nu}\eta_{\rho\sigma}/2}{2(x^2)^4},
 \quad I_{\mu\nu}=\eta_{\mu\nu}-\frac{2x_\mu x_\nu}{x^2}.
\]
Their overall constants include the prescribed spectral boundary
normalization. Let $\Pi\subset\cH$ be the Hilbert sum of these wave
channels, omitting zero-norm channels. Kernel covariance defines an
auxiliary unitary conformal action on $\Pi$: transform tube point vectors
with their indicated tensor matrices and weights, observe preservation
of their scalar products, and extend from their dense span. This is a
comparison representation, not yet a symmetry of the physical theory.
Its real localization spaces can equivalently be constructed as those
of the corresponding positive generalized free fields.

The following rest axes are adapted to the fixed null direction:
\begin{equation}\label{F:tetrad}
 e_z^{(n)}=\frac{M}{b}n-\frac pM,\qquad
 e_i^{(n)}=e_i+\frac{k_i}{b}n.
\end{equation}
They are orthogonal to $p$, have squared norm $-1$, and are mutually
orthogonal. Define the constant spin labels $j_A$ and $s_{AB}$ by the
components of $J/(iM)$ and $t/M^2$ in this tetrad, and put $s_z=s_{zz}$.
Then
\begin{equation}\label{F:longitudinal-labels}
 V_v=ibc_V,\qquad T_{vv}=b^2c_T,\qquad
 c_V=l+j_z,\qquad c_T=l/3+s_z.
\end{equation}
The phases refer to real derivatives of Hermitian fields. Expansion of
\eqref{F:tetrad} shows that every stress component is a linear combination
of constant labels with polynomial coefficients in $a,b,k,M$, allowing
only powers of $b$ in the denominators. Its norm is bounded by a constant
times $w_\alpha w_\beta$, where
\begin{equation}\label{F:component-weights}
 w_u=a,\qquad w_i=\sqrt{ab},\qquad w_v=b.
\end{equation}
For example $M,|k|\le\sqrt{2ab}$, and the apparent $1/b$ factors in the
rest axes acquire a compensating $M$ or $k$ from the rest amplitude.

\subsection{Scalar comparison copies}
\label{F:scalar-comparisons}

For each label in the following table, some tangential derivative of
$M^\nu d_*$ is the vacuum wave of a physical differential field:
\begin{equation}\label{F:comparison-table}
 \begin{array}{c|c}
 d_*&\nu\\\hline
 c_V,c_T&0\\
 l,j_z,s_z&2\\
 j_i,s_{zi}&1\\
 (s_{ij})_{\perp\mathrm{TF}}&2
 \end{array}
\end{equation}
Here the derivative has multiplier $b^m$ for some integer $m\ge0$,
up to a fixed phase. To see the physical content of this claim, set
$q_i=be_i+k_i n$, a vector whose contractions are differential
combinations with polynomial coefficients. The identities
\[
 M^2l=-\Theta,\qquad
 bM^2j_z=M^2V_v/i+b\Theta,\qquad
 b^2M^2s_z=M^2T_{vv}+b^2\Theta/3
\]
prove the scalar and longitudinal entries. Since $p\cdot q_i=0$,
\[
 V(q_i)=ibMj_i,\qquad T(n,q_i)=b^2Ms_{zi},\qquad
 \bigl(T(q_i,q_j)\bigr)_{\perp\mathrm{TF}}
       =b^2M^2(s_{ij})_{\perp\mathrm{TF}}.
\]
The scalar improvement in \eqref{F:wave-decomposition} drops from the
first mixed tensor contraction and from the transverse traceless part.
Finally $c_V,c_T$ come directly from \eqref{F:longitudinal-labels}.
All factors $M^2,b,k_i$ in these identities are polynomial momentum
multipliers; no inverse differential operator has been declared local.
The scalar reality phase may be chosen as $-i$ on the $j_i,s_{zi}$
copies and as $1$ on the other entries. This accounts for the factor
$i$ in a real tangential derivative.

Suppose a region $R$ is invariant under positive
$v$ translations, as are the cuts used below. Derivatives in $v$ can be
removed in the real closure of its wave space. Choose real
$\rho\in C_c^\infty(0,1)$ with integral one, and set
$\rho_R(v)=R^{-1}\rho(v/R)$. For a test function $f$ supported in the
region, the wave of
$(1-\rho_R*_{v})^m f$ has multiplier
$(1-\widehat\rho(Rb))^m$ and converges to that of $f$ by dominated
convergence, since $b>0$ almost everywhere. Each such test has its first
$m$ $v$-moments zero. Its $m$-fold primitive is therefore compactly
supported; choosing the primitive to vanish at $v=-\infty$ keeps its
support in the positive-$v$ hull of the original support and its positive
translates. Thus it is the $m$th derivative of a real test in the same
region. With the fixed derivative phase, the scalar generalized free
wave space of dimension $2+\nu$ on each label embeds in $X(R)$.

For $W=\{u<0,v>0\}$ this copy has precisely its own scalar wedge
Tomita data. Indeed, its embedding intertwines boosts and the wedge
reflection, including the derivative phase; hence its projection reduces
$J_We^{-\pi B}$, and the fixed real space is the scalar one. This does
not say that its curved-cut real space exhausts physical localization.

In light-front variables the measure is, up to one common fixed
normalization,
\begin{equation}\label{F:lightfront-measure}
 \dd^4p=\frac{\dd b}{b}\,\dd^2k\,M\dd M,
 \qquad b>0, M>0.
\end{equation}
For a scalar copy of dimension $2+\nu$, put $Y=i\partial_k$ and define
\begin{equation}\label{F:Hankel-generator}
 H_\nu=b(|Y|^2+h_\nu),\qquad
 h_\nu=-\partial_M^2-M^{-1}\partial_M+\nu^2/M^2.
\end{equation}
The operator $h_\nu$ is the regular nonnegative Hankel realization on
$L^2(\R_+,M\dd M)$: the Hankel transform with kernel $J_\nu(Mr)$
conjugates it to multiplication by $r^2$. In particular, the logarithmic
solution at zero is excluded when $\nu=0$.

The associated null conformal map is
\begin{equation}\label{F:null-conformal-map}
 (u,y,v)\longmapsto
 (\gamma u,\gamma y,v+c\gamma|y|^2),\qquad
 \gamma=(1-2cu)^{-1},\qquad c>0.
\end{equation}
It maps $W$ diffeomorphically to
\[
 -\frac1{2c}<u<0,\qquad
 v>\frac{|y|^2}{2u+1/c}.
\]
Two-point conformal covariance therefore maps the scalar wedge real
space onto the scalar real space of this region. Its unitary implementer
is $e^{icH_\nu}$. For verification, apply \eqref{F:Hankel-generator} to
$M^\nu e^{ip\cdot x}$. The direct expression is
\[
 H_\nu(M^\nu e^{ip\cdot x})
 =\left(2u^2a-2u k\cdot y+b|y|^2-2iu(\nu+2)\right)
       M^\nu e^{ip\cdot x}.
\]
For complex tube coordinates, the squares in this formula are analytic
bilinear squares. Multiplication by $i$ gives the infinitesimal
coordinate action
\[
 2u^2\partial_u+2uy_i\partial_i+|y|^2\partial_v
       +2u(2+\nu),
\]
which integrates to \eqref{F:null-conformal-map} with the scalar weight.
This calculation is valid on tube point vectors: their imaginary
future displacement gives exponential damping in $b$ and Gaussian
damping in $k,M$ at fixed $b$; at $M=0$ their radial behavior is
$M^\nu$ times an even smooth function, the regular Hankel boundary
condition. Applying any power of $b(|Y|^2+h_\nu)$ retains an integrable
polynomial times this damping. The factors of $b$ cancel the inverse
Gaussian variances as $b\downarrow0$. Consequently differentiation
occurs in the operator domain. Uniqueness of the unitary evolution
identifies this realization with the kernel-covariance implementer.

\subsection{The physical null-cut generators}
\label{F:cuts}

For a real transverse profile $f$, define
\begin{equation}\label{F:cut-definition}
 W_f=\operatorname{int}\left\{u<0:
 v+\frac{|y-z|^2}{-2u}\ge f(z)\text{ for every }z\in\R^2\right\}.
\end{equation}
Equivalently, put
\[
 \mathcal L_d f(y)=\sup_z\left(f(z)-\frac{|y-z|^2}{2d}\right);
 \qquad W_f=\{u<0, v>\mathcal L_{-u}f(y)\}^{\circ}.
\]
The allowed profiles initially are nonnegative bounded smooth functions
and nonnegative circular parabolas of positive curvature. We subsequently
allow constants of either sign, and sums of a bounded smooth function
with a positive-curvature circular parabola. All bounded smooth profiles
used in estimates have the stated derivatives bounded. Write
$p_d(y)=|y-d|^2$ and $W=W_0$.

For a nonnegative profile, $W_f\subset W$ has nonempty open interior and
an open spacelike complement. Reeh--Schlieder and the test-space density
above make the vacuum vectors of finite factorized polynomials from one
small double cone dense. Their affiliated realizations have bounded polar
cutoffs in the algebra, so $\Omega$ is cyclic for $\cM(W_f)$. Locality and
cyclicity in the complement make it separating.
The boost sends $W_f$ to $W_{e^sf}$, which is contained in $W_f$ for
$s\ge0$. Every hypothesis of \eqref{F:HSMI} is therefore met, and it
defines $H_f\ge0$ with
\begin{equation}\label{F:cut-modular}
 \begin{gathered}
 S(W_f)=B-H_f,\qquad
 \cM(W_f)=\Ad(e^{iH_f})\cM(W),\\
 e^{isS(W_f)}=e^{i(1-e^s)H_f}e^{isB},\qquad
 H_{\lambda f}=\lambda H_f\quad(\lambda>0).
 \end{gathered}
\end{equation}
The last assertion follows by conjugating the modular difference by a
boost and applying its affine covariance.

Translations $e^{icb}$ act inward on both $W$ and $W_f$ for $c\ge0$.
Equation~\eqref{F:Borchers} says that both modular groups scale $b$ by
the same factor; their quotient in \eqref{F:cut-modular} consequently
commutes with the whole $b$-group. Hence $H_f$ and $b$ strongly commute.
A constant translation gives
\begin{equation}\label{F:constant-shifts}
 H_{f+c}=H_f+cb.
\end{equation}
Use this identity to define signed profiles bounded below; it agrees
with their modular difference and is independent of the positive shift
chosen. Inclusion and modular-log order then give, for bounded $g$,
\begin{equation}\label{F:bounded-cut-bound}
 -\norm{g}_\infty b\le H_g\le\norm{g}_\infty b.
\end{equation}
The inequalities can first be tested on vectors with $b$ in a compact
band, smooth for the relevant affine-group representations. On such a
band they assert a bounded self-adjoint extension, and exhaustion of
the bands gives the displayed decomposable-operator inequality.

\Needspace{6\baselineskip}
\begin{proposition}[Parabolic generators]\label{F:parabolic-generator}
On $\{b>0\}$ there is a boost-independent internal Hilbert space
$\mathcal K$ and a nonnegative self-adjoint operator $h$ on it such that
\begin{equation}\label{F:parabolic-form}
 \cH_{b>0}\simeq
 L^2(\R_+,\dd b/b)\otimes L^2(\R^2,\dd^2k)\otimes\mathcal K,
 \qquad H_{p_d}=b\bigl(|Y-d|^2+h\bigr),\quad Y=i\partial_k.
\end{equation}
Here $G=bY$, $P_i=-k_i$, and
$(e^{isB}\psi)(b,k)=\psi(e^sb,k)$.
The operators $H_{p_d}$ vanish on $\ker b$. Moreover
$J_WbJ_W=b$, $J_WYJ_W=Y$, and $J_WhJ_W=h$.
\end{proposition}

\begin{proof}
Put $H=H_{p_0}$. Completion of the square in
\eqref{F:cut-definition} shows that $W_{p_0/(2R)}$ is the portion,
with $u<0$, of the future cone with tip $(u,v,y)=(-R,0,0)$.
The null rotations about this tip have generators $G_i+RP_i$.
They preserve the algebra and the vacuum, so commute with its modular
group
$e^{i(1-e^s)H/(2R)}e^{isB}$.
Set $\tau=(1-e^s)/(2R)$. Boost conjugation of the rotation group then
gives, at the level of Weyl elements, the transformation
\[
 ((1-2R\tau)\beta,R\beta)\longmapsto(\beta,R\beta)
\]
for coefficients of $(G,P_\perp)$; its central coefficient $b$ is fixed.
For a fixed sufficiently small nonzero $\tau$, vary $R$ over an open
positive interval. Two distinct values give linearly independent
coefficient pairs. Multiplying the corresponding Weyl elements, with
the same central Weyl phase on both sides, yields the full shear
\begin{equation}\label{F:parabolic-shear}
 \Ad(e^{i\tau H})G_i=G_i,\qquad
 \Ad(e^{i\tau H})P_i=P_i+2\tau G_i.
\end{equation}
The group law extends it to every real $\tau$.

Positive energy confines the nonzero momentum spectrum on $\ker b$ to
one null ray. A small spatial rotation takes this ray to a disjoint
ray; a vector supported there would be orthogonal to its arbitrarily
small rotations, contradicting strong continuity. Thus $\ker b$ is
precisely the zero-momentum subspace. On it $P_i=0$, and
\eqref{F:parabolic-shear} forces $G_i=0$. Repeating in all null directions
makes every connected Lorentz generator vanish there. In particular
$B=0$, and $e^{isB}He^{-isB}=e^sH$ forces $H=0$ on this subspace.

For $b>0$, the null rotations and transverse translations form the Weyl
pair $G_i=bY_i$, $P_i=-k_i$. Their regular representation is the
Schr\"odinger representation with a multiplicity space; the boost action
on $\log b$ trivializes that multiplicity over $b$. This gives the
representation in \eqref{F:parabolic-form}. The operator $b|Y|^2$
implements exactly the shear \eqref{F:parabolic-shear}. Consequently
$e^{-itb|Y|^2}e^{itH}$ commutes with the transverse Weyl representation
and with $b$. It is a one-parameter unitary group on the multiplicity
space, fiber by fiber. Boost homogeneity identifies its generator as
$bh$ with $h$ independent of $b$. Since $H\ge0$, its expression
$b(|Y|^2+h)$ implies $h\ge0$: localize the $Y$ variable near zero and
test an arbitrary spectral interval of $h$.

Transverse translation covariance replaces $Y$ by $Y-d$, proving the
formula for $p_d$. The modular-conjugation relation for $H$ and $b$,
together with geometric conjugation of the null rotations, gives the
three statements involving $J_W$.
\end{proof}

The use of a regular Weyl representation in this proof is the
multiplicity form of the Stone--von Neumann theorem, or equivalently
translation imprimitivity. No assertion that the internal space consists
only of one-field waves is made. In the channels of $\Pi$, $M$ and the
spin labels are internal variables.

\Needspace{6\baselineskip}
\begin{proposition}[Bounded cuts commute with the parabolic data]
\label{F:bounded-commutation}
For bounded smooth $g$,
\[
 H_g=bh_g,\qquad \norm{h_g}\le\norm{g}_\infty,
\]
where $h_g$ is independent of $b$ and strongly commutes with $Y,h$.
\end{proposition}

\begin{proof}
The formula \eqref{F:parabolic-form} and constant shifts make all
parabolic generators commute. Consequently
$\Ad(e^{itH_{p_d}})\cM(W_q)=\cM(W_{q+tp_d})$ for every
positive-curvature circular parabola $q$ and $t\ge0$.
To pass to a bounded $g$ we need an algebra-generation argument.
For $x_0=(u_0,v_0,y_0)$, $u_0<0$, its null shadow is
\[
 q_{x_0}(z)=v_0+\frac{|z-y_0|^2}{-2u_0},\qquad
 W_{q_{x_0}}=I^+(x_0)\cap\{u<0\}.
\]
Given a double cone compactly contained in $W_g$, choose $x_0$ slightly
timelike before its lower tip, still inside $W_g$. Its shadow strictly
dominates $g$ and its future cone contains the double cone. Conversely,
$q\ge g$ implies $W_q\subset W_g$. Inner generation therefore proves
\begin{equation}\label{F:parabolic-generation}
 \cM(W_g)=\bigvee_{q>g}\cM(W_q).
\end{equation}
For $W_{g+tp_d}$ the interior-shadow construction supplies a generating
subfamily of dominating parabolas of curvature strictly greater than
$t$. Indeed lower curvature cannot dominate $g+tp_d$ at transverse
infinity, while equality corresponds to the boundary value
$u_0=-1/(2t)$ and cannot be the shadow of an interior point. Thus the
join defining $\cM(W_{g+tp_d})$ may be restricted to all dominating
parabolas of curvature greater than $t$. Subtraction of $tp_d$ bijects
this restricted family with the positive-curvature parabolas dominating
$g$. Arbitrary dominating parabolas may have curvature exactly $t$;
they are not needed for the generating family. Taking joins yields
\[
 \Ad(e^{itH_{p_d}})\cM(W_g)=\cM(W_{g+tp_d}).
\]
Conjugate the modular group and use
$\Ad(e^{itH_{p_d}})B=B-tH_{p_d}$. Applying the same identity to
$\lambda g$ and $\lambda t$, then dividing by $\lambda$ by boost
homogeneity, gives
\begin{equation}\label{F:parabola-addition}
 H_{g+tp_d}=tH_{p_d}
    +e^{i\lambda tH_{p_d}}H_g e^{-i\lambda tH_{p_d}}
 \quad(\lambda>0).
\end{equation}
There is no need to differentiate an unspecified unbounded sum. On a
band $\epsilon\le b\le L$, $H_g$ is bounded by
$L\norm{g}_\infty$. The second term in \eqref{F:parabola-addition} is
independent of $\lambda$, and its strong limit as $\lambda\downarrow0$
is $H_g$. It follows that $H_g$ commutes with the full group of every
$H_{p_d}$. Exhausting the $b$ bands proves strong commutation globally.
In the joint calculus of the commuting parabolic family,
\[
 H_{p_d}-H_{p_0}=b(-2d\cdot Y+|d|^2).
\]
These differences recover $Y$ on $b>0$, and $H_{p_0}/b-|Y|^2$ recovers
$h$. Hence $H_g$ commutes with $Y,h$. Finally
\eqref{F:bounded-cut-bound} gives the bounded operator $h_g=H_g/b$,
and boost homogeneity makes it independent of $b$.
\end{proof}

We next justify the reflected localization needed for lower bounds.
The wedge conjugation reflects $(u,v,y)$ to $(-u,-v,y)$. Even if its
action on an arbitrarily chosen bounded double-cone algebra has not
been postulated, it produces a net compatible with the original one
in spacelike separation. Indeed every bounded operation in a double
cone belongs to all containing wedge algebras; their images under
$J_W$ are the corresponding reflected, twist-corrected wedge algebras
by \eqref{F:BW-sandwich} and Poincar\'e covariance. Two compact
spacelike separated double cones admit separating complementary
wedges. To see this, write their causal tips as $a_i,b_i$. A future
null linear functional separating $b_1$ from $a_2$, and a second one
separating $b_2$ from $a_1$, define the two wedge boundaries; their
strict inequalities place the cones in opposite wedges. Locality of
these wedge algebras gives the claim.

In particular, the reflected region $J_WW_{-g}$ is spacelike to $W_g$.
This also follows directly by minimizing the sum of the two quadratic
penalties in \eqref{F:cut-definition}; for $s,t>0$ its minimum is
$|y_--y_+|^2/(2(s+t))$. Thus
\begin{equation}\label{F:reflected-order}
 J_WX(W_{-g})\subset X(W_g)',\qquad H_g\ge-H_{-g}.
\end{equation}
For the second implication use modular-log order,
$J_WBJ_W=-B$, and $J_WH_gJ_W=H_g$. The latter follows from the
positive-cut conjugation relation and the constant-shift formula.

\Needspace{6\baselineskip}
\begin{lemma}[Taylor comparison]\label{F:taylor-bound}
If $g$ is bounded and smooth with bounded Hessian, then
\begin{equation}\label{F:taylor-equation}
 |h_g-g(Y)|\le C h,\qquad
 C=\tfrac12\sup_y\norm{g''(y)}.
\end{equation}
In particular
$\norm{(h_g-g(Y))\psi}\le C\norm{h\psi}$ for $\psi\in\Dom h$.
\end{lemma}

\begin{proof}
For every transverse center $z$,
\[
 g(y)\le g(z)+\nabla g(z)\cdot(y-z)+C|y-z|^2.
\]
The right side is a signed positive-curvature parabola if $C>0$.
Inclusion, modular-log order, and \eqref{F:parabolic-form} give
\[
 h_g\le g(z)+\nabla g(z)\cdot(Y-z)+C|Y-z|^2+Ch.
\]
Applying the same argument to $-g$ and then \eqref{F:reflected-order}
gives the lower inequality with both occurrences of $C$ negated.
For this fixed $g$, all operators involved commute strongly by
Proposition~\ref{F:bounded-commutation}. Impose the inequalities first
for a countable dense set of $z$ and then, on each joint spectral fiber,
let $z$ tend to the value of $Y$. They become
$-Ch\le h_g-g(Y)\le Ch$ in a commuting spectral calculus, which proves
both assertions. If $C=0$, boundedness makes $g$ constant and the
result is \eqref{F:constant-shifts}.
\end{proof}

\Needspace{6\baselineskip}
\begin{proposition}[Scalar spectral comparison]\label{F:spectral-comparison}
On any scalar copy from \eqref{F:comparison-table}, with its embedding
into the ambient internal space understood, one has
\begin{equation}\label{F:spectral-equation}
 \mathbf1_{(\lambda,\infty)}(h)
       \mathbf1_{[0,\lambda]}(h_\nu)=0,\qquad
 \norm{h^s\psi}\le\norm{h_\nu^s\psi}\quad(s>0).
\end{equation}
The norm assertion includes the corresponding inclusion of domains.
It does not assert that $h$ leaves the scalar copy invariant.
For bounded profiles $f\le g$ one also has
\begin{equation}\label{F:bounded-spectral-order}
 \mathbf1_{(\lambda,\infty)}(h_f)
       \mathbf1_{(-\infty,\lambda]}(h_g)=0.
\end{equation}
\end{proposition}

\begin{proof}
Let $j_\nu$ embed the scalar comparison copy. The derivative-removal
argument in Section~\ref{F:scalar-comparisons} gives its real wave
localization inside the physical cut space. By
\eqref{F:null-conformal-map} and \eqref{F:cut-modular},
\[
 j_\nu e^{icH_\nu}X_\nu(W)
       \subset e^{icH_{p_0}}X(W)\qquad(c>0).
\]
The embedding intertwines the flat wedge boosts, so
Lemma~\ref{F:strip-comparison} applies. It also intertwines the
external variables $b,Y$. To see the latter point independently of
the internal operator $h$, a null rotation $N_\beta$ fixes $n,b,M$
and sends $k$ to $k+b\beta$ and $e_i$ to $e_i+\beta_i n$.
The adapted axes in \eqref{F:tetrad} therefore satisfy
$N_\beta e_A^{(n)}(p)=e_A^{(n)}(N_\beta p)$ for $A=z,1,2$.
The constant labels and their fixed derivative phases transform
accordingly. Thus $j_\nu$ intertwines the null-rotation groups
generated by $G=bY$, as well as $b$; disintegration on $b>0$
gives the claimed $Y$ intertwining. This uses the Poincar\'e action
on the comparison waves and does not require $h$ to preserve their
individual channels. Both sides are decomposable in
$b,Y$, and $j_\nu$ intertwines these variables. In the common spectral
fibers their generators are
$b(|Y|^2+h_\nu)$ and $b(|Y|^2+h)$. Subtracting the common scalar
$b|Y|^2$ and dividing by $b>0$ in the band-inclusion conclusion gives
the first assertion in \eqref{F:spectral-equation}. More formally,
apply the whole-generator conclusion for rational band endpoints,
localize $(b,Y)$ to shrinking measurable rectangles, and use monotone
convergence of spectral projections. The power bound follows by
integrating spectral tails as in Lemma~\ref{F:strip-comparison}.

If $f\le g$, the cut inclusion is $W_g\subset W_f$. Apply the same
lemma with source $g$, target $f$, and the identity embedding. If needed
add the same constant to make both profiles nonnegative. After
disintegrating $b$, removing this common constant gives
\eqref{F:bounded-spectral-order}. This is spectral order obtained from
an analytic strip, rather than from ordinary quadratic-form order.
\end{proof}

\Needspace{8\baselineskip}
\subsection{Normal displacement and the required form domain}

\begin{lemma}[Normal form bound]\label{F:lipschitz-form}
Let $g$ be bounded and smooth with Lipschitz constant $C_1$, and write
$H_g(-\delta)=e^{-i\delta a}H_g e^{i\delta a}$. For $\delta,\lambda>0$,
\begin{equation}\label{F:normal-displacement}
 \pm\bigl(H_g-H_g(-\delta)\bigr)
 \le\delta\left(\frac a\lambda+\frac{\lambda C_1^2}2 b\right)
\end{equation}
as forms on a joint momentum-smooth core. The commutator form with
$a$ extends continuously to $\Dom\sqrt{1+a+b}$. If
$\psi\in\bigcap_{m\ge1}\Dom(1+a+b)^m$, then
$H_g\psi\in\Dom\sqrt a$.
\end{lemma}

\begin{proof}
The quadratic sup-convolution obeys
$\mathcal L_{d+\delta}f=\mathcal L_d(\mathcal L_\delta f)$.
Completing the square in the Lipschitz estimate gives
\[
 \mathcal L_\delta(\lambda g)(y)
 \le\lambda g(y)+\frac{\delta\lambda^2 C_1^2}{2}.
\]
Hence the cut with profile
$\lambda g+\delta\lambda^2C_1^2/2$, translated by $u=-\delta$, lies
inside $W_{\lambda g}$. The translated boost is $B-\delta a$.
Modular-log order thus gives
\[
 B-\delta a-\lambda H_g(-\delta)
     -\frac{\delta\lambda^2C_1^2}{2}b
 \le B-\lambda H_g.
\]
This is the upper sign of \eqref{F:normal-displacement}. Apply $J_W$,
which fixes $a,b,H_g$ and reverses the translation parameter by
antiunitarity, and then translate by $-\delta$; this gives the lower
sign. One may initially use vectors smooth for translations and the
boost, then remove the boost smoothing in the resulting form bound.

Put $E=1+a+b$ and
$Q_\lambda=a/\lambda+\lambda C_1^2 b/2$. Differentiating at
$\delta=0$ on momentum-smooth vectors defines the Hermitian form
\[
 c_g(\phi,\psi)=i\bigl(
  \ip{a\phi}{H_g\psi}-\ip{H_g\phi}{a\psi}\bigr).
\]
Differentiation is legitimate using $H_g=bh_g$, bounded $h_g$, and
translation smoothness of both vectors; it does not presume that
$H_g$ preserves $\Dom a$. The two inequalities give
$|c_g(\psi,\psi)|\le\ip{\psi}{Q_\lambda\psi}$. Polarization, or
representation of this Hermitian form after $Q_\lambda^{1/2}$, gives
\begin{equation}\label{F:commutator-form}
 |c_g(\phi,\psi)|\le
 \norm{Q_\lambda^{1/2}\phi}\,
 \norm{Q_\lambda^{1/2}\psi}.
\end{equation}
For fixed $\lambda$ this proves continuity in the $E^{1/2}$ form norm.

For the domain assertion set $\chi=H_g\psi$. The $b$ bound and strong
$b$ commutation imply $\chi\in\Dom\sqrt{1+b}$. On a joint spectral
core for $a,b$,
\[
 \ip{a\phi}{\chi}
   =\ip{H_g\phi}{a\psi}-i c_g(\phi,\psi),\qquad
 |\ip{H_g\phi}{a\psi}|
 \le\norm{g}_\infty\norm{\sqrt b\,\phi}
                           \norm{\sqrt b\,a\psi}.
\]
Both terms are continuous in $\norm{\sqrt E\,\phi}$. Adding
$\ip{(1+b)\phi}{\chi}$ consequently gives
\[
 |\ip{E\phi}{\chi}|\le C_\psi\norm{\sqrt E\,\phi}.
\]
The spectral theorem characterizes this inequality as
$\chi\in\Dom\sqrt E$: test with spectral truncations of $\chi$ and
let their endpoints tend to infinity. Thus $H_g\psi\in\Dom\sqrt a$,
as required.
\end{proof}

When opposite localization of conjugated bounded operations is needed,
one may therefore use the reflected, twist-corrected net just described.
It is spacelike compatible with the original net. Conjugating a chosen
original affiliated realization by $J_W$ gives an affiliated realization
in the reflected algebra, with its polynomial restriction on
$J_W\mathcal D_{\mathrm{old}}$. On generated factorized vacuum words,
its vacuum and adjoint-vacuum data are the reflected-adjoint word data
from the wedge boundary calculation; admitted compact joint word data follow by the graph limits
above. No equality of $J_W\mathcal D_{\mathrm{old}}$ with
$\mathcal D_{\mathrm{old}}$ is used. The construction can also be repeated with past cuts
of the original net. Section~\ref{sec:null-current} will identify their
bounded-profile generators from the physical current and thereby compare
the resulting moments.
